% !PRINTED 613
\weberpage{613}
%<<CONT
werden können. Nehmen wir als Gruppe die Ordnung $O' = [Q]$ mit dem Führer $Q$, die durch die Kongruenzbedingung definiert ist, daß jede Zahl in $O'$ nach dem Modul $Q$ mit einer rationalen Zahl kongruent sein soll, so haben wir eine Kongruenzgruppe, deren Exkludent und Modul $= Q$ ist. Die Klassen entsprechen den Formenklassen der Diskriminante $D = Q^2\Delta$. Der zweite Fall schließt den ersten in sich (für $Q = 1$). Der Klassenkörper ist hier der Klassenkörper $\mathfrak{K}(D)$, den wir zum Unterschied von allgemeineren als \gesperrt{Ordnungskörper} bezeichnen wollen.

Denn ist $(k)$ eine Klasseninvariante, so ist (\S\,122)
\[
(k)^p \equiv (k) \quad (\mathrm{mod}\ \mathfrak{P})
\]
nur dann befriedigt, wenn $p$ durch die Hauptklasse der Diskriminante $D$ darstellbar ist. Dann zerfällt $p$ in $\Omega$ in zwei konjugierte Primideale $\mathfrak{p}$, $\mathfrak{p}'$, die durch \gesperrt{Zahlen} in $O$ oder in $O'$ darstellbar, also Hauptideale sind.

Der Körper $\mathfrak{K}(\Delta)$ ist der \gesperrt{Klassenkörper von $\Omega$ im engeren Sinne}, der den Idealklassen dieses Körpers entspricht, und auch als \gesperrt{Haupt-Klassenkörper} $\mathfrak{K}(\Omega)$ bezeichnet sein mag.

In jeder Idealklasse, sowohl nach $O$ als nach $O'$, gibt es also unendlich viele Primideale ersten Grades. Mit anderen Worten: Ist $\mathfrak{a}$ ein beliebiges (zu $Q$ teilerfremdes) Ideal, so gibt es unendlich viele Primideale $\mathfrak{p}$, die mit $\mathfrak{a}$ nach $O'$ äquivalent sind. Dann sind nach \S\,95 auch die den Idealen $\mathfrak{a}$ und $\mathfrak{p}$ entsprechenden Formen äquivalent, und $p$ ist durch die zu $\mathfrak{a}$ gehörige \gesperrt{Form darstellbar}. Damit ist der Satz bewiesen:

1. \gesperrt{Durch jede primitive Form der Diskriminante $D$ sind unendlich viele Primzahlen darstellbar}\footnote{Dieser Satz, sowohl für negative als für positive Diskriminanten, ist zuerst von \gesperrt{Dirichlet} bewiesen, der Beweis aber nur in einem speziellen Fall publiziert (Bericht der Berliner Akademie 1840, Werke Bd.\ I, S.\ 497). Der ausgeführte Beweis ist von H. \gesperrt{Weber} gegeben (Mathematische Annalen XX, 1882). Eine gleichfalls von \gesperrt{Dirichlet} herrührende Verallgemeinerung, auf die wir im nächsten Paragraphen zurückkommen, ist von A. \gesperrt{Meyer} bewiesen (Crelles Journ., Bd.\ 103). Der hier im Text gegebene Beweis beruht auf anderer Grundlage, bezieht sich aber freilich einstweilen (solange der Klassenkörper für positive Diskriminanten nicht bekannt ist) nur auf negative Diskriminanten.}.

\begin{center}
\S\,167. \textbf{Primzahlen in Linearformen.}
\end{center}

Wir definieren nun eine Zahlgruppe $A$ in $\Omega$ folgendermaßen:
%CONT>>
